% generated by review3.py appcode at 2026-09-23T05:25:10.885156+00:00; do not edit

\paragraph{Baseline schema} (\code{src/schemaeval/schemas/medium.py})
\begin{lstlisting}[language=Python]
import re
from typing import List, Literal

from pydantic import BaseModel, Field, model_validator

CANONICAL_ID_PATTERN = r"^[a-z][a-z0-9_]{2,30}$"
_ID_RE = re.compile(CANONICAL_ID_PATTERN)

class Entity(BaseModel):
    name: str = Field(min_length=1, max_length=100)
    entity_type: Literal["person", "place", "organization", "concept", "other"]
    canonical_id: str = Field(
        pattern=CANONICAL_ID_PATTERN,
        description="lowercase_snake_case identifier, e.g. 'manchester_united'",
    )

class IdentifyOutput(BaseModel):

    entities: List[Entity] = Field(min_length=2, max_length=5)
    question_type: Literal[
        "who", "what", "when", "where", "why", "how", "comparison"
    ]
    primary_entity_id: str = Field(
        description="Must match one of entities[].canonical_id exactly"
    )
    entity_count: int = Field(
        ge=2, le=5, description="Must equal len(entities)"
    )

    @model_validator(mode="after")
    def _consistency(self):
        ids = {e.canonical_id for e in self.entities}
        if self.primary_entity_id not in ids:
            raise ValueError(
                f"primary_entity_id '{self.primary_entity_id}' not in "
                f"entities canonical_ids {sorted(ids)}"
            )
        if self.entity_count != len(self.entities):
            raise ValueError(
                f"entity_count={self.entity_count} != "
                f"len(entities)={len(self.entities)}"
            )
        return self

class Fact(BaseModel):
    entity_id: str = Field(
        pattern=CANONICAL_ID_PATTERN,
        description="Reference to a canonical_id from the identify step",
    )
    fact: str = Field(min_length=10, max_length=400)
    relevance: float = Field(ge=0.0, le=1.0)

class GatherOutput(BaseModel):

    facts: List[Fact] = Field(min_length=2, max_length=8)
    fact_count: int = Field(
        ge=2, le=8, description="Must equal len(facts)"
    )
    coverage_notes: str = Field(min_length=20, max_length=250)

    @model_validator(mode="after")
    def _consistency(self):
        if self.fact_count != len(self.facts):
            raise ValueError(
                f"fact_count={self.fact_count} != len(facts)={len(self.facts)}"
            )
        return self

class AnswerOutput(BaseModel):

    answer: str = Field(min_length=20, max_length=400)
    supporting_entity_ids: List[str] = Field(min_length=1, max_length=5)
    confidence: float = Field(ge=0.0, le=1.0)
    confidence_bucket: Literal["low", "medium", "high"] = Field(
        description=(
            "Must match confidence numerically: "
            "confidence < 0.5 => 'low'; "
            "0.5 <= confidence <= 0.8 => 'medium'; "
            "confidence > 0.8 => 'high'"
        )
    )

    @model_validator(mode="after")
    def _consistency(self):
        expected = (
            "low" if self.confidence < 0.5
            else "high" if self.confidence > 0.8
            else "medium"
        )
        if self.confidence_bucket != expected:
            raise ValueError(
                f"confidence={self.confidence} implies bucket='{expected}', "
                f"got '{self.confidence_bucket}'"
            )
        for eid in self.supporting_entity_ids:
            if not _ID_RE.match(eid):
                raise ValueError(
                    f"supporting_entity_id '{eid}' fails canonical_id pattern "
                    f"'{CANONICAL_ID_PATTERN}'"
                )
        return self
\end{lstlisting}

\paragraph{Prompt variants} (\code{src/schemaeval/prompts.py})
\begin{lstlisting}[language=Python]
import json

_BACKTICK_FENCE = "`" * 3

def _baseline(schema_json_str: str) -> str:
    return (
        "You must respond with a single JSON object matching this schema.\n"
        "Return ONLY the JSON — no markdown code fences, no prose, no explanation.\n\n"
        f"Schema:\n{schema_json_str}"
    )

def _aggressive(schema_json_str: str) -> str:
    return (
        "CRITICAL: You MUST respond with a single JSON object matching the "
        "schema below.\n"
        "Do NOT use markdown formatting.\n"
        f"Do NOT wrap your response in code fences ({_BACKTICK_FENCE}json "
        f"or {_BACKTICK_FENCE}).\n"
        "Do NOT include any prose, explanation, or commentary.\n"
        "Your response MUST begin with { as the very first character.\n"
        "Any response starting with backticks or prose will be rejected.\n\n"
        f"Schema:\n{schema_json_str}"
    )

def _few_shot(schema_json_str: str) -> str:
    return (
        "You must respond with a single JSON object matching this schema.\n"
        "Return ONLY the JSON — no markdown code fences, no prose, no explanation.\n\n"
        "Example of a CORRECT response format:\n"
        '{"example_field": "example_value", "count": 3}\n\n'
        "Example of an INCORRECT response format (do not do this):\n"
        f"{_BACKTICK_FENCE}json\n"
        '{"example_field": "example_value", "count": 3}\n'
        f"{_BACKTICK_FENCE}\n\n"
        "The second example is wrong because it uses markdown code fences. "
        "Do not wrap JSON in backticks or code blocks.\n\n"
        f"Schema:\n{schema_json_str}"
    )

_PROMPT_BUILDERS = {
    "baseline": _baseline,
    "aggressive": _aggressive,
    "few_shot": _few_shot,
}

def build_system_prompt(schema: type, variant: str = "baseline") -> str:
    if variant not in _PROMPT_BUILDERS:
        raise ValueError(
            f"Unknown prompt variant '{variant}'. "
            f"Available: {sorted(_PROMPT_BUILDERS.keys())}"
        )
    schema_json_str = json.dumps(schema.model_json_schema(), indent=2)
    return _PROMPT_BUILDERS[variant](schema_json_str)

def available_variants() -> list[str]:
    return sorted(_PROMPT_BUILDERS.keys())
\end{lstlisting}

\paragraph{medium\_v2 changes}
\begin{lstlisting}
--- src/schemaeval/schemas/medium.py
+++ src/schemaeval/schemas/medium_v2.py
@@ -16,5 +16,4 @@
 
 class IdentifyOutput(BaseModel):
-
     entities: List[Entity] = Field(min_length=2, max_length=5)
     question_type: Literal[
@@ -48,14 +47,13 @@
         description="Reference to a canonical_id from the identify step",
     )
-    fact: str = Field(min_length=10, max_length=400)
+    fact: str = Field(min_length=10, max_length=800)      # was 400
     relevance: float = Field(ge=0.0, le=1.0)
 
 class GatherOutput(BaseModel):
-
     facts: List[Fact] = Field(min_length=2, max_length=8)
     fact_count: int = Field(
         ge=2, le=8, description="Must equal len(facts)"
     )
-    coverage_notes: str = Field(min_length=20, max_length=250)
+    coverage_notes: str = Field(min_length=20, max_length=1000)  # was 250
 
     @model_validator(mode="after")
@@ -68,6 +66,5 @@
 
 class AnswerOutput(BaseModel):
-
-    answer: str = Field(min_length=20, max_length=400)
+    answer: str = Field(min_length=20, max_length=1200)   # was 400
     supporting_entity_ids: List[str] = Field(min_length=1, max_length=5)
     confidence: float = Field(ge=0.0, le=1.0)
\end{lstlisting}

\paragraph{medium\_v3 changes}
\begin{lstlisting}
--- src/schemaeval/schemas/medium.py
+++ src/schemaeval/schemas/medium_v3_loose_regex.py
@@ -4,5 +4,5 @@
 from pydantic import BaseModel, Field, model_validator
 
-CANONICAL_ID_PATTERN = r"^[a-z][a-z0-9_]{2,30}$"
+CANONICAL_ID_PATTERN = r"^[a-zA-Z0-9][a-zA-Z0-9_\-]{2,30}$"
 _ID_RE = re.compile(CANONICAL_ID_PATTERN)
 
@@ -12,5 +12,5 @@
     canonical_id: str = Field(
         pattern=CANONICAL_ID_PATTERN,
-        description="lowercase_snake_case identifier, e.g. 'manchester_united'",
+        description="alphanumeric identifier, may include hyphens/underscores",
     )
 
\end{lstlisting}
